% ============================================================
% PRM393 - Section 8: Conclusion & Future Work (Bilingual)
% ============================================================
\section{\bi{Kết luận và hướng phát triển}{Conclusion \& Future Work}}
\label{sec:concl}

\subsection{\bi{Kết luận}{Conclusion}}

\bi{
Báo cáo này đã trình bày toàn diện kiến trúc kỹ thuật và kết quả hiện thực của
ứng dụng di động OmniCast Mobile trong khuôn khổ môn học PRM393. Bốn mục tiêu
kỹ thuật đặt ra ban đầu đều được chứng minh bằng mã nguồn và dữ liệu thực
nghiệm:
}{
This paper has presented the comprehensive software architecture and
empirical findings of OmniCast Mobile for the PRM393 course. The four original
engineering goals have been validated through concrete source code and
rigorous benchmarks:
}

\begin{enumerate}[leftmargin=*,itemsep=2pt]
  \item \bi{Hiện thực thành công kiến trúc Clean Architecture 4 tầng kết hợp
  12 module BLoC/Cubit, duy trì sự phân tách rõ ràng giữa giao diện hiển thị
  và logic nghiệp vụ.}{%
  A 4-tier Clean Architecture coupled with 12 decoupled BLoC/Cubit units,
  strictly segregating reactive UI rendering from domain calculations.}

  \item \bi{Xây dựng động cơ ngoại tuyến hoàn chỉnh trên SQLite với cơ chế
  phát lại đột biến (\texttt{pending\_ops}) và nhắc lịch cục bộ trước 15 phút,
  bảo đảm tính khả dụng 100\% khi không có mạng.}{%
  An offline-first persistence engine driven by SQLite, mutation replay
  queuing (\texttt{pending\_ops}), and exact 15-minute advance local alarms.}

  \item \bi{Tích hợp trình phát đa năng \texttt{OmniPlayer} hỗ trợ phát trực
  tiếp WebRTC/HLS độ trễ thấp từ MediaMTX, đồng thời duy trì kênh tương tác hai
  chiều (Live Chat, Floating Reactions, Live Donations).}{%
  The multi-engine \texttt{OmniPlayer} enabling ultra-low-latency WebRTC/HLS
  live streams, backed by bi-directional WebSocket overlays for chat,
  reactions, and donations.}

  \item \bi{Đạt độ tin cậy tuyệt đối với 91 ca kiểm thử tự động đạt tỷ lệ vượt
  qua 100\% trong quy trình CI/CD GitHub Actions.}{%
  Absolute software reliability substantiated by 91 automated unit and widget
  tests achieving a 100\% pass rate on GitHub Actions CI runners.}
\end{enumerate}

\subsection{\bi{Đánh giá hiệu năng và hạn chế đã biết}{Performance Evaluation \& Limitations}}

\bi{
Thử nghiệm trên các thiết bị Android và iOS thực tế cho thấy ứng dụng duy trì
tỷ lệ khung hình ổn định ở mức 58--60 FPS khi cuộn danh sách EPG và hiển thị
các hiệu ứng hoạt họa nổi. Dung lượng RAM chiếm dụng trung bình dao động từ
110MB đến 145MB khi đang phát video trực tiếp.

Tuy nhiên, ứng dụng vẫn còn hai hạn chế kỹ thuật:
}{
Field profiling across physical Android and iOS devices indicates steady
rendering performance at 58--60 FPS during intensive EPG scrolling and floating
reaction animations. Memory usage averages between 110MB and 145MB during
continuous live stream decoding.

Nevertheless, two known technical limitations remain:
}

\begin{itemize}[leftmargin=*,itemsep=2pt]
  \item \bi{Chưa hỗ trợ giải mã phần cứng nâng cao đối với codec AV1 trên các
  dòng máy Android đời cũ, dẫn đến việc phải chuyển đổi dự phòng sang định dạng
  HLS H.264.}{%
  Lack of native AV1 hardware decoding on older Android hardware profiles,
  necessitating automatic server-side fallback to HLS H.264 profiles.}
  \item \bi{Cơ chế đồng bộ ngoại tuyến hiện áp dụng chiến lược \emph{Last-Write-Wins}
  đơn giản cho danh sách theo dõi, chưa triển khai cấu trúc dữ liệu CRDT để
  xử lý xung đột đồng thời trên nhiều thiết bị.}{%
  Offline synchronization employs a pragmatic \emph{Last-Write-Wins}
  policy rather than Conflict-Free Replicated Data Types (CRDTs), which may
  yield minor discrepancies during simultaneous cross-device updates.}
\end{itemize}

\subsection{\bi{Hướng phát triển tiếp theo}{Future Work}}

\bi{
Các định hướng nâng cấp trong giai đoạn tiếp theo bao gồm:
}{
Roadmap horizons for upcoming releases include:
}

\begin{enumerate}[leftmargin=*,itemsep=2pt]
  \item \bi{Tích hợp chế độ phát hình trong hình (Picture-in-Picture --- PiP)
  và phát âm thanh trong nền (Background Audio Playback) cho các luồng tin tức
  và podcast.}{%
  Integrating native Picture-in-Picture (PiP) and background audio services
  for broadcast news and podcast feeds.}
  \item \bi{Mở rộng hỗ trợ nền tảng màn hình lớn Android TV và Apple TV với
  giao diện điều khiển từ xa D-Pad.}{%
  Expanding support for Android TV and Apple TV form factors with D-Pad
  remote control navigation focus states.}
  \item \bi{Ứng dụng mạng nơ-ron học máy cục bộ (TensorFlow Lite) trên thiết
  bị để đề xuất chương trình cá nhân hóa mà không cần gửi dữ liệu xem về máy
  chủ trung tâm.}{%
  Deploying lightweight on-device recommendation models via TensorFlow Lite
  to compute privacy-preserving suggestions locally.}
\end{enumerate}
